\lecture{10}{Sep 15 10:15}{Gaussian Elimination}

\section{Gaussian Elimination with Back Substitution}

\subsection{Coefficient Matrix}
\begin{definition}[Matrix]\label{def:matrix}
    A \textbf{matrix} is a rectangular table of scalars from a field.

    The number of rows and columns is called the \textbf{order} of the matrix. A matrix of order $m \times n$ matrix has $m$ rows and $n$ columns. Individual scalar values are called the \textbf{entries} of the matrix.
\end{definition}
\begin{notation}
    We usually name matrices with uppercase Roman letters and matrix elements have the corresponding lowercase letter.
    For example:
    \[
        A = \begin{bmatrix}
            a_{11} & a_{12} & a_{13} \\
            a_{21} & a_{22} & a_{23} \\
            a_{31} & a_{32} & a_{33}
        \end{bmatrix}
    \]
\end{notation}
\begin{definition*}[Augmented Coefficient Matrix]\label{def:augmented-coefficient-matrix}
    Any linear system of $m$ equations with $n$ unknowns is represented fully by an $m \times \left(n+1\right)$ matrix, called the \textbf{augmented coefficient matrix}:
    \[
        \begin{bmatrix}
            a_{11} & a_{12} & \cdots & a_{1n} & b_1 \\
            a_{21} & a_{22} & \cdots & a_{2n} & b_2 \\
            \vdots & \vdots &  \vdots  & \vdots & \vdots \\
            a_{m1} & a_{m2} & \cdots & a_{mn} & b_m
        \end{bmatrix}
    \]
\end{definition*}

\begin{definition}[Row Equivalent]\label{def:row-equivalent}
    Let $A$ and $B$ be matrices of the same order. $A$ is \textbf{row equivalent} to $B$ if we can change $A$ into $B$ by a finite sequence of elementary row operations.
\end{definition}
\begin{note}
    Row equivalence is a valid equivalence relation because it is reflexive, symmetric, and transitive.
\end{note}
\begin{definition}[Pivot]\label{def:pivot}
    A matrix row containing all zeros is called a \textbf{null row}. In any row that is not a null row, the leftmost non-zero entry is called the \textbf{leading coefficient} or the \textbf{pivot}
\end{definition}
\subsection{Row Echelon Form}
\begin{definition}[REF]\label{def:ref}
    A matrix is said to be in \textbf{row echelon form} (REF) if:
    \begin{enumerate}
        \item In every row that is not a null row the leading coefficient is to the right of all the leading coefficients of the rows above it (if any)
        \item All null rows, if any exist, are below all non-null rows
    \end{enumerate}
\end{definition}
\begin{definition}[eliminated system]\label{def:eliminated-system}
    A linear system whose augmented coefficient matrix is in row echelon form
\end{definition}
\begin{definition}[Restricted and Free Variables]\label{def:restricted-and-free-variables}
    In an eliminated linear system, $x_j$ is called a \textbf{restricted variable} if there is a pivot in the $j$-th column of its coefficient matrix. Otherwise, it is called a \textbf{free variable}.
\end{definition}
\begin{note}
    In a solution set of an eliminated system, the restricted variables are specified in terms of the free variables. Each free variable in the general solution $n$-tuple is represented by a free parameter.
\end{note}

\subsection{Gaussian Elimination}
Given a matrix order of $m \times n$
\begin{enumerate}
    \item If all entries are zero, the matrix is already in REF. Stop
    \item If the matrix has only one row, it is already in REF. Stop
    \item Identify the first (leftmost) column containing a non-zero entry, column $j$
    \item If the first (topmost) entry in that column is zero, swap with column with non-zero entry in column $j$
    \item If there are any non-zero entries left in column $j$ below the topmost row, add multiples of the topmost row to the rows of those entries to make them zero
    \item Ignore the topmost row and repeat setps $1-5$ on the remaining sub-matrix
\end{enumerate}
\begin{note}
    This algorithm works on any field
\end{note}